\section{Discussion}
The experiments support a bounded conclusion: ContextTrace provides a
reproducible local evidence-forensics contract and strong pre-review
failure-detection behavior on the frozen samples. They do not establish that it
is universally more accurate than RAG evaluators, nor that its suggested fixes
improve human debugging. The matched RAGAS comparison favors different systems
on different common metrics, and unsupported diagnostic outputs are not a fair
basis for dominance claims.

The ablations suggest that diagnostic modules matter jointly. Source state and
citation alignment substantially reduce false greens, while abstention logic is
important for root-cause assignment. This supports exposing the chain rather
than only its final score. At the same time, the simulated RQ4 result shows that
more structure can reduce perceived actionability: a detailed report may be
correct yet cognitively expensive or insufficiently prioritized. Human study
design should therefore measure time and preference, not only diagnosis match.

\paragraph{Claim boundary.}
We claim a benchmarked, local-first evidence-chain forensics system with frozen
pre-review results. We do not claim broad SOTA, independent review, human
actionability improvement, or real-world truth verification. Dataset-specific
metrics must retain sample, baseline, and review status. The machine-readable
gate currently blocks broad SOTA while review is pending.

\paragraph{Reproducibility.}
The full workflow pins case IDs, profiles, bootstrap seed, candidate files,
dependency versions, command, hashes, runtime, and cost metadata.
Table~\ref{tab:table6-reproducibility} summarizes the run. A post-simulated-review
rerun from a committed revision preserves all case hashes and reported metrics,
applies zero corrections, and includes sensitivity artifacts separately.

% Generated by paper/generate_tables.py; do not edit manually.
\begin{table*}[t]
\centering
\scriptsize
\resizebox{\textwidth}{!}{%
\begin{tabular}{llll}
\toprule
 Item & Review status & Value & Status \\
\midrule
 Source revision & pre\_review\_paper\_candidate & 15cf09f9f9d7bfdd6285dae2603edad12d8a3ff9 & frozen \\
 Command & pre\_review\_paper\_candidate & python benchmarks/contexttrace\_bench/reproduce\_arr\_tables.py --full --output-dir 'benchmarks\textbackslash{}contexttrace\_bench\textbackslash{}out\textbackslash{}arr\_full' & full \\
 Python & pre\_review\_paper\_candidate & 3.8.10 & recorded \\
 Runtime seconds & pre\_review\_paper\_candidate & 249.847 & recorded \\
 Bootstrap & pre\_review\_paper\_candidate & 400 / 20260612 & frozen \\
 ContextTrace API cost & pre\_review\_paper\_candidate & 0.0 & local \\
 Candidate cost & same-ID cached baseline & N/A & cached/unreported \\
 Paper result eligible & pending independent review & False & pending review \\
\bottomrule
\end{tabular}%
}
\caption{Table6 Reproducibility. Review status and unavailable measurements are shown explicitly.}
\label{tab:table6-reproducibility}
\end{table*}

